\section{Introduction}

Software systems that combine large language models (LLMs) with tools, retrieval, and multi-step control flow are often described as \emph{agentic}: they repeatedly observe an environment, update internal state, and emit actions until a termination condition holds \citep{russell2020aima}. Surveys of ``agentic AI'' emphasize planning depth, tool diversity, and memory mechanisms, but rarely tie these informal axes to the classical theory of computation.

The Chomsky hierarchy partitions formal languages by generative power and, equivalently, by the automaton class needed to recognize them \citep{hopcroft2006automata}. Type~3 (regular) languages correspond to finite-state control; Type~2 (context-free) to a single stack; Type~1 (context-sensitive) to linear-bounded workspace; Type~0 to unrestricted Turing machines.

This paper does \emph{not} assert that a deployed LLM instantiates a particular grammar on natural language. Instead, we treat an \emph{agentic program} as a specified controller---often a finite template of prompts, parsers, and API calls---whose \emph{interaction traces} form a formal language over finite grammars of observable events (finite sets of observation and action terminals after discretization). Restrictions on that controller's state space and memory discipline yield a hierarchy of trace languages that parallels Types~3--0. The contribution is definitional: we supply a common formal basis for comparing agent designs, benchmarking, and peer review, in line with separating claims from evidence in multi-model evaluation \citep{react2023}.

\section*{Claims and scope}

We state comparative claims only at the level of \emph{agent schemata} (Section~\ref{sec:schemata}) and their idealized trace languages. Empirical measurements of real LLM agents must map implementations to these schemata; Section~\ref{sec:eval} outlines how to pre-register such mappings.
